% =====================================================================
%  GMN3 §4.1--4.3  形式线算子及其非凡的穿墙性质
%  原文：Lamination-Paper.tex 2021--2606 行；PDF p.28--38（印刷 p.27--37）
%  编号公式 (4.1)--(4.29) 共 29 条，全部在本块内（块内引用一律 \eqref）；
%  图 9 幅（4.1--4.9），图名照抄原文 \insfig 第一参数。
%  跨块引用（§3.4、§4.4、§5.3、§12、式 3.29/3.38/3.48/3.49）暂写原文编号并加 % TODO。
% =====================================================================

% —— 原文宏 \ch / \hW / \qbinom 在 preamble.tex 中尚未定义，此处临时补上；
%    合稿时应移入 preamble.tex，并删除本块这三行（\providecommand 不会覆盖已有定义）。

\section{形式线算子及其非凡的穿墙性质}\label{sec:Formal-line-operators}

halo 穿墙图像与强正性猜想结合起来，威力十分强大。强正性不给自旋特征标中的相消留下任何余地：$y^m X_\gamma$ 的系数是希尔伯特空间的真实维数，并按红外电荷与 $so(3)$ 自旋分次。当我们改变 $\zeta$、改变理论的参数与真空期望值时，halo 穿墙会具体地增添或删去这些希尔伯特空间的子空间，但维数必须始终非负。这是一个非常强的约束；在本小节中，我们想就几个具体的 $\CN=2$ 理论探讨它的推论，这些理论具有简单且已被充分理解的「原味」BPS 谱。

\subsection{腔与生成函数}\label{subsec:Chambers}

在本小节中，我们作一些形式定义，用以概括生成函数 (3.38)% TODO \eqref{eq:form-gen}
的行为。

引入非交换环 $\IZ[y,y^{-1},X_\gamma]/I$，其中理想 $I$ 由关系 $yX_\gamma = X_\gamma y$ 与 $X_\gamma X_{\gamma'} = y^{\langle\gamma,\gamma'\rangle}X_{\gamma+\gamma'}$ 生成。现在回顾式 (3.49)% TODO \eqref{eq:Chamber-Space}
所定义的空间 $\Xi$。正如那里所说，$\Xi$ 被划分成若干腔。我们用 $c$ 标记这些腔，其中 $c$ 取遍某个指标集。我们定义\emph{强正形式线算子}为一族元素 $F(c)\in\IZ[y,y^{-1},X_\gamma]/I$，使得

\begin{enumerate}

\item 在腔 $c^+,c^-$ 之间的墙 $\hW(\gamma)$ 两侧，
\begin{equation}
F(c^+) = S_\gamma F(c^-) S_\gamma^{-1},
\end{equation}
其中 $c^+$、$c^-$ 分别是 $\Im(Z_\gamma/\zeta)$ 取正值、取负值的腔，$S_\gamma$ 定义于式 (3.48)% TODO \eqref{eq:Spin-KS-ii}
。

\item 在 $\Xi$ 的每个腔 $c$ 中，
\begin{equation}\label{eq:finite-positive-sum}
F(c) = \sum_\gamma P^c_\gamma(y) X_\gamma
\end{equation}
其中 $P^c_\gamma(y)$ 是 $SU(2)$ 的某个（真实的、而非虚拟的）表示的特征标。

\end{enumerate}

上面第二项定义了「强正」条件。我们注意到，我们给出的定义与 Fock 和 Goncharov 在研究簇系综（cluster ensemble）时引入的\emph{泛洛朗多项式}（universal Laurent polynomial）\cite{MR2567745} 密切相关。

若强正性猜想成立，则线算子的生成函数 $F(L)$ (3.38)% TODO \eqref{eq:form-gen}
是一个强正形式线算子。我们把形式线算子记作 $F$。与真实的线算子一样，在考虑线算子族时，$F$ 会承受紫外参数中的单值变换。当然，如果我们在 $\CB\times\IC^*$ 上、而非在其万有覆盖上定义腔，那么在这个空间里也会有单值变换。

在接下来的几个小节中，我们将考察一些简单 $\CN=2$ 理论中形式线算子的例子。在这些情形中，我们会发现腔结构有一些重要的简化。例如，我们将只考虑只涉及单个超多重态的墙。此时第 3.4 节% TODO \ref{subsec:Halo-Wall-Crossing}
的全套量子双对数技术有点小题大做——我们只需直接应用式 (3.29)% TODO \eqref{eq:transform}
，按需移除或添加 halo。单个超多重态的穿墙规则——\emph{从} ${\rm Im}(Z_{\gamma_h}/\zeta)<0$ 的区域\emph{到} ${\rm Im}(Z_{\gamma_h}/\zeta)>0$ 的区域——算出来简单地就是
\begin{equation}\label{eq:explicit-dilog-conj}
X_{\gamma_c} \to \begin{cases}
\sum_{j=0}^N \ch\left(\Lambda^j \rho_N\right) X_{\gamma_c + j\gamma_h} & \langle\gamma_h,\gamma_c\rangle = N > 0 \\[2pt]
\sum_{j=0}^\infty (-1)^j \ch\left(S^j \rho_N\right) X_{\gamma_c + j\gamma_h} & \langle\gamma_h,\gamma_c\rangle = -N < 0
\end{cases}
\end{equation}
这里 $\rho_N$ 是 $SU(2)$ 的 $N$ 维表示。式 \eqref{eq:explicit-dilog-conj} 的第一行显然保持正性。然而第二行表明，要维持正性是高度非平凡的。注意在第一种情形中，内积的符号使我们进入一个支持核心 $\gamma_c$ 周围 $\gamma_h$ 粒子 halo 的区域；而在第二种情形中，我们进入的区域不支持这样的 halo 构型。\emph{从} ${\rm Im}(Z_{\gamma_h}/\zeta)>0$ 的区域\emph{到} ${\rm Im}(Z_{\gamma_h}/\zeta)<0$ 的区域的变换规则同样算得
\begin{equation}\label{eq:explicit-dilog-conj-ii}
X_{\gamma_c} \to \begin{cases}
\sum_{j=0}^\infty (-1)^j \ch\left(S^j \rho_N\right) X_{\gamma_c + j\gamma_h} & \langle\gamma_h,\gamma_c\rangle = N > 0 \\[2pt]
\sum_{j=0}^N \ch\left(\Lambda^j \rho_N\right) X_{\gamma_c + j\gamma_h} & \langle\gamma_h,\gamma_c\rangle = -N < 0
\end{cases}
\end{equation}
同样，当我们进入支持 halo 构型的区域时，正性明显地保持；否则不然。当然，先应用式 \eqref{eq:explicit-dilog-conj} 再应用式 \eqref{eq:explicit-dilog-conj-ii} 必须给出恒等算子。这蕴涵一个多少非平凡的恒等式
\begin{equation}\label{eq:Char-Ident}
\sum_{j+k=\ell} (-1)^j \ch\left(S^j \rho_N\right) \ch\left(\Lambda^k \rho_N\right) = \delta_{\ell,0},
\end{equation}
对整数 $\ell\geq 0$ 成立。这就是在用式 \eqref{eq:explicit-dilog-conj} 从支持 halo 构型的区域过渡到不支持的腔时，halo 被「移除」的方式。

关于式 \eqref{eq:explicit-dilog-conj} 与式 \eqref{eq:explicit-dilog-conj-ii}，我们想作最后一点评论。注意当表达式有限时，它明显地是强正的；而当它无限时则不然。这提示：$X_\gamma$ 展开的有限性与强正性相关。在第 \ref{subsec:Formal-line-ops-U1} 节与第 \ref{sec:formaln3} 节中，我们将在一些简单场论中构造强正形式线算子。然而在第 4.4 节% TODO \ref{subsec:formal-SU2}
中，我们将处理一个更复杂的场论，其中我们难以证明所构造的形式线算子的强正性。尽管如此，受所有显式例子与上述观察的启发，我们想猜想：\emph{若一个形式线算子在单个腔中是强正的，并且在所有腔中都有有限的达布展开，那么它事实上就是强正的。}

\subsection{形式线算子与 $U(1)$ 规范理论}\label{subsec:Formal-line-ops-U1}

有一个方便的玩具模型，它抓住了形式线算子行为的若干关键方面：带单个电荷为 $1$ 的超多重态的 $U(1)$ 规范理论。这个理论在紫外没有良好定义，但它确实抓住了良定义理论在库仑分支 $\CB$ 的奇点附近的行为。

在此例中，$\CB^*$ 是去心圆盘，$\Gamma$ 是秩二的局部系统。$\Gamma$ 的纤维由 $\gamma_1,\gamma_2$ 生成，$\langle\gamma_1,\gamma_2\rangle = +1$，绕原点的逆时针单值变换作用为
\begin{equation}
\gamma_1 \to \gamma_1+\gamma_2, \qquad \gamma_2 \to \gamma_2.
\end{equation}
我们把 $p\gamma_1+q\gamma_2$ 记作 $(p,q)$。

我们可以取 $Z_{\gamma_2}(u)=u$。我们不需要 $Z_{\gamma_1}(u)$ 的显式公式，因为非零的「原味」BPS 简并度只有 $\Omega(\pm\gamma_2)=+1$。「原味」BPS 简并度不发生穿墙。

我们预期能够找到对应于各种电荷与磁荷的 Wilson--'t Hooft 线算子的线算子。由于带电粒子的存在带来单圈 beta 函数，我们预期会看到反常 R 对称性的效应：一次 R 对称变换要求 $\theta$ 角平移，而后者又通过 Witten 效应引起带磁荷对象电荷的平移。这导致线算子标签在 $\zeta$ 旋转 $2\pi$ 之下发生单值变换。这种单值变换与穿墙之间会有漂亮的相互作用。

万有覆盖 $\widehat\CB \times \widehat\IC^*$ 就是 $\IC\times\IC$，坐标为 $(\log u,\log\zeta)$。给定我们的「原味」BPS 谱，墙是
\begin{equation}
\hW(\gamma_2) = \amalg_{n\in\IZ}\left\{(\log u,\log\zeta): \arg u-\arg\zeta = (2n+1)\pi\right\},
\end{equation}
\begin{equation}
\hW(-\gamma_2) = \amalg_{n\in\IZ}\left\{(\log u,\log\zeta): \arg u-\arg\zeta = 2n\pi\right\},
\end{equation}
从而腔可以取为
\begin{equation}
c_n := \left\{(\log u,\log\zeta): n\pi < \arg u-\arg\zeta < (n+1)\pi\right\}.
\end{equation}

让我们现在尝试构造形式线算子 $F^{(n)}_{p,q}$：规定它在腔 $c_n$ 中的值为 $F^{(n)}_{p,q}(c_n) = X_{p,q}$。这就是电荷为 $q$、磁荷为 $p$ 的 Wilson--'t Hooft 线算子的朴素生成函数。变换法则 \eqref{eq:explicit-dilog-conj} 与 \eqref{eq:explicit-dilog-conj-ii} 随即确定 $F^{(n)}_{p,q}$ 在所有其他腔中的取值，于是我们可以问：它们满足强正性约束吗？

首先，若 $p=0$ 则没有变换，故对所有腔 $c$ 都有 $F^{(n)}_{0,q}(c)=X_{0,q}$。这当然满足形式正性的约束。由于它们实际上与 $n$ 无关，我们把这些形式线算子简单记作 $F_{0,q}$。它们是电荷为 $q$ 的 Wilson 线算子，没有有趣的穿墙，因为「原味」粒子只带电荷。

现在注意，在稳定性条件 $\langle\gamma_h,\gamma_c\rangle\,{\rm Im}Z_{\gamma_h}/\zeta > 0$ 中，我们只需要考虑 $\gamma_h=\pm\gamma_2$，而且事实上对这两个电荷该条件相同。若 $\gamma_c=(p,q)$，则稳定性条件变为 $-p\,{\rm Im}Z_{\gamma_2}/\zeta>0$。因此当 $p>0$ 时，$c_n$ 在 $n$ 为奇数时可以有 halo 构型、在 $n$ 为偶数时不能；而当 $p<0$ 时，$c_n$ 在 $n$ 为偶数时可以有 halo 构型、在 $n$ 为奇数时不能。见图 \ref{fig:chamber-u1-even} 与图 \ref{fig:chamber-u1-odd}。

\insfig{chamber-u1-even}{最简单的腔与强正形式线算子：带单个 BPS 超多重态的 $U(1)$ 规范理论，此处 $n=0\bmod 2$、$p>0$。对 halo 形成稳定的区域用阴影标出。}
\insfig{chamber-u1-odd}{最简单的腔与强正形式线算子：带单个 BPS 超多重态的 $U(1)$ 规范理论，此处 $n=1\bmod 2$、$p<0$。对 halo 形成稳定的区域用阴影标出。}

在分析 $F^{(n)}_{p,q}$ 的行为时，有四种情形要考虑：$n$ 可为偶数或奇数，$p$ 可正可负。

\begin{enumerate}

\item $n=0\bmod 2$，$p>0$。此时我们从 $c_n$ 走到 $c_{n+1}$ 会穿过类型为 $\hW(\gamma_2)$ 的墙。我们应用式 \eqref{eq:explicit-dilog-conj-ii} 的第二行，得到
\begin{equation}\label{eq:n-even-p-pos}
F^{(n)}_{p,q}(c_{n+1}) = \sum_{j=0}^p \qbinom{p}{j} X_{p,q+j}.
\end{equation}
同样，从 $c_n$ 走到 $c_{n-1}$ 我们会穿过类型为 $\hW(-\gamma_2)$ 的墙。此时应用式 \eqref{eq:explicit-dilog-conj} 的第一行，得到
\begin{equation}
F^{(n)}_{p,q}(c_{n-1}) = \sum_{j=0}^p \qbinom{p}{j} X_{p,q-j}.
\end{equation}
当我们从 $c_{n+1}$ 继续走到 $c_{n+2}$ 时，本可以应用我们的形式规则并使用恒等式 \eqref{eq:Char-Ident}，但反推回去更简单：回忆稳定性条件不允许腔 $c_{n+2}$ 中有 halo。因此，从 $c_{n+2}$ 穿过墙 $\hW(-\gamma_2)$ 到 $c_{n+1}$ 时，我们必须是在产生一个由电荷 $-\gamma_2$ 的粒子组成的 halo。这个操作确实会给出 \eqref{eq:n-even-p-pos}，当且仅当
\begin{equation}\label{eq:u1-monod}
F^{(n)}_{p,q}(c_{n+2}) = X_{p,q+p}.
\end{equation}
类似地，继续走到腔 $c_{n-2}$，我们得到 $F^{(n)}_{p,q}(c_{n-2})=X_{p,q-p}$。这个模式显然向两个方向延续下去，因此 $F^{(n)}_{p,q}$ 是一个形式线算子。此外，此式表明 $F^{(n)}_{p,q} = F^{(n+2)}_{p,q+p}$，因此不失一般性，我们可以把这些算子归约为由腔 $c_0$ 中的值确定的单一族。$n\to n+2$ 时电荷的平移有两种独立的解释：若把它解释为 $\zeta$ 变量单值变换的结果，那它就是紫外标签的单值性 $F_{p,q}\to F_{p,q+p}$；而若把它解释为 $u$ 平面中单值变换的结果，它表达的则是红外电荷格 $\Gamma$ 绕奇异轨迹 $u=0$ 的单值性。

最后，这个例子说明：在把 $F$ 按 $X_\gamma$ 展开时，代表核心的项与代表 halo 的项之间没有不变的区别。的确，在从 $c_n$ 穿过 $\hW(\gamma_2)$ 到 $c_{n+1}$ 时，核心电荷是 $(p,q)$；而当从 $c_{n+2}$ 穿过 $\hW(-\gamma_2)$ 到 $c_{n+1}$ 时，核心电荷是 $(p,q+p)$。这与我们的物理图像一致：核心与 halo 的区别，只有当我们非常靠近一道墙时才变得分明。

\item $n=0\bmod 2$，$p<0$。此时稳定性条件允许在腔 $c_n$ 中有 halo、而在 $c_{n\pm1}$ 中不允许，所以我们预期会有麻烦，因为当我们试图继续到相邻的腔 $c_{n\pm1}$ 时不能有 halo 构型。的确，此时穿过墙 $\hW(\gamma)$ 进入腔 $c_{n+1}$ 必须应用式 \eqref{eq:explicit-dilog-conj-ii} 的第 1 行。用这个变换作用于 $X_{p,q}$，显然给出违反强正性的表达式。因此 $F^{(n)}_{p,q}$ 不是形式线算子。

\item $n=1\bmod 2$，$p>0$。这与 $(n=0\bmod 2,\,p<0)$ 的情形类似，不给出形式线算子。

\item $n=1\bmod 2$，$p<0$。这与 $(n=0\bmod 2,\,p>0)$ 的情形类似。它确实给出形式线算子，其中
\begin{equation}\label{eq:n-odd-p-neg}
F^{(n)}_{p,q}(c_{n+1}) = \sum_{j=0}^{|p|} \qbinom{|p|}{j} X_{p,q-j},
\end{equation}
且 $F^{(n)}_{p,q}(c_{n+2})=X_{p,q+p}$。同样有 $F^{(n)}_{p,q}=F^{(n+2)}_{p,q+p}$，故不失一般性可以只考虑 $F^{(1)}_{p,q}$。

\end{enumerate}

总之，由于 $F^{(n)}_{p,q}=F^{(n+2)}_{p,q+p}$，为了构造那些在某个腔中约化为单项式的强正形式线算子，只需考虑腔 $n=0$ 与 $n=1$。所得的算子与整数对 $(p,q)$ 一一对应，因此我们可以直接把它们记作 $F_{p,q}$。若 $p=0$，则 $F_{0,q}=X_{0,q}$ 在所有腔中成立。若 $p>0$，则 $F_{p,q}=X_{p,q}$ 在腔 $c_0$ 中成立；若 $p<0$，则 $F_{p,q}=X_{p,q}$ 在腔 $c_1$ 中成立。这些显然是单（simple）线算子，而且事实上它们是唯一的单线算子：给定任何形式线算子，我们可以把它拆成磁荷为正的部分、不带磁荷的部分与只带负磁荷的部分。显然，磁荷为正的部分只需在某个 $n$ 为偶数的腔中分解成单项式，就能分解为单的构件。对磁荷为负的部分（取 $n$ 为奇数）以及对不带磁荷的部分（任意 $n$），类似的说法也成立。形式线算子 $F_{p,q}$ 应当与 $U(1)$ 理论中 't Hooft--Wilson 圈的保护自旋特征标等同起来。

即使在这个最简单的例子里，形式线算子的代数也已经有些非平凡，虽然它还简单到可以显式写出。的确，$F_{\pm1,0}$、$F_{0,\pm1}$ 已经生成一个环。人们容易手算出关系 $F_{1,0}F_{-1,0}=1+yF_{0,1}$，但更高阶的乘法涉及对 $q$-二项式系数的繁琐操作。环关系可以如下容易地算出。首先，很容易验证
\begin{equation}
F_{p,q}F_{r,s} = y^{ps-qr}F_{p+r,q+s}
\end{equation}
在 $p\geq 0,r\geq 0$ 或 $p\leq 0,r\leq 0$ 时成立，因为此时可以在算子为单项式的腔 $c_0$ 或 $c_1$ 中求值。$p,r$ 异号的情形要非平凡得多。为确定起见，设 $p\geq 0$、$r\leq 0$。那么按定义：
\begin{equation}
F_{p,q}F_{r,s} = \begin{cases}
X_{p,q}\left(\sum_{j=0}^{|r|}\qbinom{|r|}{j} X_{r,s+j}\right) & {\rm in}~~c_0 \\[2pt]
\left(\sum_{j=0}^p \qbinom{p}{j} X_{p,q+j}\right) X_{r,s} & {\rm in}~~c_1
\end{cases}
\end{equation}
于是执行乘法即得：
\begin{equation}
F_{p,q}F_{r,s} = \begin{cases}
y^{ps-qr}\sum_{j=0}^{|r|}\qbinom{|r|}{j} y^{pj} X_{p+r,q+s+j} & {\rm in}~~c_0 \\[2pt]
y^{ps-qr}\sum_{j=0}^p \qbinom{p}{j} y^{-rj} X_{p+r,q+s+j} & {\rm in}~~c_1
\end{cases}
\end{equation}
现在，要辨认右端对应于哪个线算子，我们应当把结果放在「磁荷为 $p+r$ 的线算子为单项式」的那个腔中来看。于是得到线算子上的乘法：
\begin{equation}\label{eq:U1-ring-rel}
F_{p,q}F_{r,s} = \begin{cases}
y^{ps-qr}\sum_{j=0}^{|r|}\qbinom{|r|}{j} y^{pj} F_{p+r,q+s+j} & p+r \geq 0 \\[2pt]
y^{ps-qr}\sum_{j=0}^p \qbinom{p}{j} y^{-rj} F_{p+r,q+s+j} & p+r \leq 0
\end{cases}
\end{equation}
我们可以容易地得到 $p\leq 0$、$r\geq 0$ 的另一情形：回忆 $F(L_1)F(L_2)$ 等于把 $F(L_2)F(L_1)$ 作替换 $y\to 1/y$。对 \eqref{eq:U1-ring-rel} 的暴力验算蕴涵关于 $y$-二项式的若干非平凡恒等式（$y$-形变的 Pascal 关系），
\begin{equation}
y^j \qbinom{p}{j} = \qbinom{p-1}{j} + y^p \qbinom{p-1}{j-1}.
\end{equation}

\subsubsection*{更高电荷}

此构造的一个自然变形，是考虑与单个电荷为 $Q$ 的超多重态耦合的 $U(1)$ 理论。也就是说，我们现在考虑 BPS 谱
\begin{equation}\label{eq:Q-BPS-spec}
\Omega(\gamma) = \begin{cases} 1 & \gamma = \pm Q\gamma_2, \\ 0 & {\rm else}. \end{cases}
\end{equation}
我们可以通过一个小技巧立刻得到形式线算子：若把理论的电荷格与磁荷格按相反的因子 $Q$ 重新标度，我们就回到 $Q=1$ 的情形。上面的代数关系适用于 $F'_{p',q'}$，其中 $F'_{p',q'} := F_{Q^{-1}p',\,Qq'}$。名副其实的线算子必须带整数个电荷。于是我们现在得到关系
\begin{equation}\label{eq:Q-U1-algebra}
F_{1,0}F_{-1,0} = \sum_{j=0}^Q \qbinom{Q}{j} y^{Qj} F_{0,Qj}.
\end{equation}

\subsection{形式线算子与 $N=3$ Argyres--Douglas 理论}\label{sec:formaln3}

我们现在描述一个带穿墙的理论的最简单例子中的形式线算子。这就是 $N=3$ Argyres--Douglas 理论。关于这个理论的更多细节，见文献 \cite{Gaiotto:2009hg} 第 9.4 节。简言之，$\CB = \IC-\{u_+,u_-\}$。万有覆盖已经相当复杂，所以我们将在 $u$ 平面上的确定区域工作，并考虑（$\zeta$ 平面万有覆盖中的）腔。对局部系统最好的描述是 $\Gamma = H_1(\Sigma_u,\IZ)^-$，其中 $\Sigma_u$ 是由
\begin{equation}
\lambda^2 = (z^3-3\Lambda^2 z+u)(dz)^2
\end{equation}
定义的、$\CB$ 上一族黎曼面的纤维，其中 $z\in\IC$。纤维 $\Gamma_u$ 是秩二的，我们取生成元 $\gamma_1,\gamma_2$ 使 $\langle\gamma_1,\gamma_2\rangle=1$。向量 $p\gamma_1+q\gamma_2$ 有时记作 $(p,q)$。在 $u=u_\pm=\pm 2\Lambda^3$ 处有两个奇点，$\Gamma$ 绕它们有单值变换。（见文献 \cite{Gaiotto:2009hg} 的式 (9.30)。）BPS 态的谱已在文献 \cite{Gaiotto:2009hg} 中确定。底空间 $\CB$ 应划分为三个区域 $\CB_s,\CB_w^\pm$，在其上 $\Gamma$ 可以平凡化。$\CB_s$ 是 $u=0$ 附近的区域，$\CB_w^\pm$ 是 $u=\infty$ 附近的两个区域。有一条穿过 $u_\pm$ 的连通墙，把 $\CB_s$ 与 $\CB_w^\pm$ 分开。我们有
\begin{equation}
\Omega(\gamma;u) = \begin{cases} 1 & \gamma = \pm\gamma_1,\pm\gamma_2 \\ 0 & {\rm else} \end{cases}
\end{equation}
对 $u\in\CB_s$，而穿墙随后确定
\begin{equation}
\Omega(\gamma;u) = \begin{cases} 1 & \gamma = \pm\gamma_1,\pm\gamma_2,\pm(\gamma_1+\gamma_2) \\ 0 & {\rm else} \end{cases}
\end{equation}
对 $u\in\CB_w^+$。$\CB_w^-$ 中的谱可由单值变换或穿墙得到。

\insfig[0.42\textwidth]{ADN3-FL-1}{$N=3$ Argyres--Douglas 理论在 $u=0$、$\Lambda>0$ 时的 BPS 墙投影到 $\zeta$ 平面上。}

在 $u=0$ 处对 $\Sigma_u$ 的 cycle 作适当选取，可使 $Z_{\gamma_1} = -i\Lambda^{5/2}K$、$Z_{\gamma_2} = -K\Lambda^{5/2}$，其中 $K>0$ 是一个（可计算的）常数。$u=0$ 处 BPS 墙到 $\zeta$ 平面的投影于是如图 \ref{fig:ADN3-FL-1} 所示。注意，当我们沿\emph{顺时针}方向行进时，穿过每道墙 $\hW(\gamma_h)$ 都是从 ${\rm Im}(Z_{\gamma_h}/\zeta)>0$ 的区域走到 ${\rm Im}(Z_{\gamma_h}/\zeta)<0$ 的区域。因此沿顺时针方向穿墙时应使用规则 \eqref{eq:explicit-dilog-conj-ii}，而沿逆时针方向穿墙时应使用规则 \eqref{eq:explicit-dilog-conj}。

让我们沿用第 \ref{subsec:Formal-line-ops-U1} 节的策略，尝试找出在 $\zeta$ 平面某个腔中约化为单项式 $X_{p,q}$ 的形式线算子。引入腔 $c_n$，其中 $\frac{\pi}{2}(1-n) < \arg\zeta < \frac{\pi}{2}(2-n)$。第一个障碍来自顺时针或逆时针方向的单次穿墙。这导致如下结果：$X_{p,q}$ 只有在下列情形下才能满足形式正性：

\begin{enumerate}

\item $p\geq 0,\ q\geq 0$，在 $n=1\bmod 4$ 的 $c_n$ 中

\item $p\leq 0,\ q\geq 0$，在 $n=2\bmod 4$ 的 $c_n$ 中

\item $p\leq 0,\ q\leq 0$，在 $n=3\bmod 4$ 的 $c_n$ 中

\item $p\geq 0,\ q\leq 0$，在 $n=0\bmod 4$ 的 $c_n$ 中

\end{enumerate}

我们现在将证明这是唯一的障碍。

\insfig[0.42\textwidth]{ADN3-FL-2}{$5$ 重覆盖的第一叶包含四个腔。}

\insfig[0.42\textwidth]{ADN3-FL-3}{沿顺时针方向继续解析延拓到第二叶。}

\insfig[0.42\textwidth]{ADN3-FL-4}{沿顺时针方向继续解析延拓到第三叶。}

\insfig[0.42\textwidth]{ADN3-FL-5}{沿顺时针方向继续解析延拓到第四叶。}

\insfig[0.42\textwidth]{ADN3-FL-6}{沿顺时针方向继续解析延拓到第五叶。从腔 $c_{20}$ 出发的顺时针穿墙产生 $X_{\gamma_1}$，于是整个过程重复。因此在万有覆盖上，腔 $21$ 与腔 $1$ 等同；在 $5$ 重覆盖上工作时，腔编号应按模 $20$ 等同。}

为确定起见，让我们从万有覆盖 $\widehat{\IC^*}$ 中满足 $0<\arg\zeta<\frac{\pi}{2}$ 的腔里的 $X_{\gamma_1}$ 出发。顺时针方向的穿墙产生图 \ref{fig:ADN3-FL-2} 所示的表达式。我们观察到一些在 $U(1)$ 例子中见过的同样现象。从 $c_1$ 到 $c_2$，我们获得一个核心电荷为 $\gamma_1$ 的 halo。从 $c_2$ 到 $c_3$，我们获得另一个 halo，但这次核心电荷是 $\gamma_1+\gamma_2$、halo 电荷是 $-\gamma_1$。从 $c_3$ 到 $c_4$，我们失去一个 halo 粒子，核心电荷仍解释为 $\gamma_1+\gamma_2$，但 halo 粒子的电荷为 $-\gamma_2$。现在，若从这个腔继续沿顺时针方向走，我们会生成一个新的表达式 $X_{\gamma_2}$。我们又从一个核心为 $\gamma_2$ 的 halo 中移走了一个 halo 粒子。不出所料，我们发现存在单值变换，因此确实必须在万有覆盖 $\widehat{\IC^*}$ 上讨论穿墙。这样做时，我们得到示意的模式：
\begin{equation}
\cdots \to X \to X \to X+X \to X+X+X \to X+X \to X \to X \to X+X \to \cdots
\end{equation}
结果表明，这个模式以 $20=4\times 5$ 为周期重复。用这个周期性对万有覆盖作商，我们得到 $\zeta$ 平面的一个 $5$ 重覆盖。每一叶包含四个腔，共 $20$ 个腔，标记为 $c_1,\dots,c_{20}$。形式线算子在这全部 $20$ 个腔中的显式表达式，由图 \ref{fig:ADN3-FL-2} 至图 \ref{fig:ADN3-FL-6} 给出。

\insfig[0.42\textwidth]{ADN3-FL-weak}{某个弱耦合区域中的腔结构。}

可惜，这还没有完全构造出强正线算子，因为我们必须检查它在弱耦合区域中是强正的。穿过边缘稳定墙进入区域 $\CB_w^+$ 时，$\gamma_1$ 与 $\gamma_2$ 的墙合并并产生第三道墙，导致 $\zeta$ 平面中出现图 \ref{fig:ADN3-FL-weak} 所示的新腔结构。每个奇数腔都分裂成两个，于是每一叶上有 $6$ 个腔，共 $30$ 个腔。注意偶数编号腔中的表达式不可能从强耦合区域发生改变。很容易验证：穿墙在奇数腔 $c_{2n+1,\pm}$ 中产生强正的元素。

于是我们已完全构造出一个强正形式线算子。利用单值变换操作，我们事实上可以得到 $5$ 个不同的形式线算子，它们由下式定义
\begin{equation}
\begin{split}
F_1(c_1) & = X_{\gamma_1} \\
F_2(c_1) & = X_{\gamma_2} \\
F_3(c_1) & = X_{-\gamma_1} + X_{\gamma_2-\gamma_1} \\
F_4(c_1) & = X_{-\gamma_1} + X_{-\gamma_1-\gamma_2} + X_{-\gamma_2} \\
F_5(c_1) & = X_{-\gamma_2} + X_{\gamma_1-\gamma_2}
\end{split}
\end{equation}

利用这些算子，我们可以如下产生若干其他强正算子。考虑 $F_{p,q}:=y^{-pq}F_1^p F_2^q$，其中 $n,m\geq 0$\zy{原文此处作「对 $n,m\geq 0$」，但式中用的是 $p,q$；应为 $p,q\geq 0$，译文照排。}。在 $u$ 处于强耦合区的腔 $c_1$ 中，它的值为 $X_{(p,q)}$。我们断言这生成一个强正形式线算子。直接检验会迅速变得非常困难，但我们可以用下述技巧来证明。首先注意，若沿顺时针或逆时针方向移动两个腔，穿墙会直截了当地产生强正的表达式。的确，我们得到显式表达式：
\begin{equation}\label{eq:Fpq-chs-N3}
\begin{split}
F_{p,q}(c_{19}) & = \sum_{j=0}^q \sum_{k=0}^{p+j} \qbinom{q}{j}\qbinom{p+j}{k} X_{p+j,q-k} \\
F_{p,q}(c_{20}) & = \sum_{j=0}^q \qbinom{q}{j} X_{p+j,q} \\
F_{p,q}(c_1) & = X_{p,q} \\
F_{p,q}(c_2) & = \sum_{j=0}^p \qbinom{p}{j} X_{p,q+j} \\
F_{p,q}(c_3) & = \sum_{j=0}^p \sum_{k=0}^{q+j} \qbinom{p}{j}\qbinom{q+j}{k} X_{p-k,q+j}
\end{split}
\end{equation}
当然，由于 $q$-二项式是真实的自旋特征标，它们的乘积也是。因此 \eqref{eq:Fpq-chs-N3} 满足强正性。现在注意，$F_1$ 在 $8$ 个腔 $c_{20},c_1$、$c_5,c_6$、$c_{10},c_{11}$、$c_{15},c_{16}$ 中成为单项式。类似地，$F_2$ 在 $8$ 个腔 $c_1,c_2$、$c_6,c_7$、$c_{11},c_{12}$、$c_{16},c_{17}$ 中成为单项式。于是在腔 $c_1,c_6,c_{11},c_{16}$ 中，表达式 $F_{p,q}$ 是单项式。具体地有
\begin{equation}\label{eq:Fpq-mon-chs}
\begin{split}
F_{p,q}(c_6) & = X_{(-q,p)} \\
F_{p,q}(c_{11}) & = X_{(-q,-p)} \\
F_{p,q}(c_{16}) & = X_{(q,-p)}
\end{split}
\end{equation}
现在，对 \eqref{eq:Fpq-mon-chs} 中的每个表达式施加与 \eqref{eq:Fpq-chs-N3} 类似的两步顺时针与逆时针变换，可以验证穿墙产生强正的表达式。这样我们就能覆盖强耦合区的全部二十个腔。第 5.2 节% TODO \ref{subsec:Seeds-from-N=2}
末尾给出的论证蕴涵：把它们延拓为弱耦合区域的强正算子没有障碍。

类似地，我们可以从其他线算子 $F_i$ 的乘积生成强正形式线算子。把上面的 $F_{p,q}$ 改记为 $F^{(1)}_{p,q}$，我们可以定义（各种情形下都有 $p\geq 0,q\geq 0$）：

\begin{enumerate}

\item $F^{(2)}_{p,q} = y^{-pq} F_2^p F_3^q$，它在 $c_2,c_7,c_{12},c_{17}$ 中是单项式

\item $F^{(3)}_{p,q} = y^{-pq} F_3^p F_4^q$，它在 $c_3,c_8,c_{13},c_{18}$ 中是单项式

\item $F^{(4)}_{p,q} = y^{-pq} F_4^p F_5^q$，它在 $c_4,c_9,c_{14},c_{19}$ 中是单项式

\item $F^{(5)}_{p,q} = y^{-pq} F_5^p F_1^q$，它在 $c_5,c_{10},c_{15},c_{20}$ 中是单项式

\end{enumerate}

我们是否应该预期 $F^{(i)}_{p,q}$ 穷尽了单的强正形式线算子之列表？我们可以勾画一个证明，它可以用第 12 节% TODO \ref{sec:Tropic}
给出的工具形式化。考虑某个形式线算子 $F$ 在腔 $c_1$ 中的展开，并看形如 $a_{n_1,n_2}X_{n_1\gamma_1+n_2\gamma_2}$ 的项，其中 $n_1\geq 0$、$n_2\geq 0$。$\tilde F = F - \sum_{n_1>0,\,n_2>0} a_{n_1,n_2}F^{(1)}_{n_1,n_2}$ 是否仍然强正？它在 $c_1$ 中是正的，并且在所有其他腔中都只有有限多项，因此它在那些腔中也必定是正的（因为从一个腔到下一个腔的穿墙只能给一个正的和增添无穷多个负项）。接下来我们可以考虑 $\tilde F$ 在腔 $c_2$ 中的展开，并减去 $F^{(2)}_{p,q}$ 的适当倍数，依此类推。过程结束时，我们得到 $F$ 分解为 $F^{(i)}_{p,q}$ 的一个正和。

现在我们可以考察这些算子生成的代数。在腔 $c_1$ 中计算，容易得到
\begin{equation}
F_1 F_3 = 1 + y F_2
\end{equation}
由于其他算子都由单值变换生成，我们有
\begin{equation}\label{eq:Deformd-N3-rels}
F_{i-1} F_{i+1} = 1 + y F_i
\end{equation}
要得到另一种顺序的乘法，取 $y\to 1/y$。
